\documentclass[aspectratio=169]{beamer}
\usetheme{metropolis}
\usepackage{booktabs}
\usepackage{listings}
\usepackage{tikz}
\usepackage{hyperref}
\usepackage{xcolor}
\usepackage{graphicx}
\usepackage{fontspec}
\usepackage{emoji}

\usetikzlibrary{shapes.geometric, arrows, positioning, chains, fit}

% Color scheme
\definecolor{darkblue}{RGB}{0,82,155}
\definecolor{lightblue}{RGB}{100,181,246}
\definecolor{darkgreen}{RGB}{46,125,50}
\definecolor{orange}{RGB}{255,152,0}
\definecolor{purple}{RGB}{156,39,176}

\setbeamercolor{frametitle}{bg=darkblue}
\setbeamercolor{progress bar}{fg=lightblue,bg=lightblue!30}

% Code listing settings
\lstset{
    basicstyle=\ttfamily\small,
    breaklines=true,
    frame=single,
    backgroundcolor=\color{gray!10},
    keywordstyle=\color{darkblue}\bfseries,
    commentstyle=\color{darkgreen}\itshape,
    stringstyle=\color{orange},
    showstringspaces=false,
    language=Python
}

\title{Lecture 5: Tokenization}
\subtitle{Week 2: From Characters to Meaningful Units}
\author{LLM Course}
\date{}

\begin{document}

% Title slide
\begin{frame}
    \titlepage
\end{frame}

% Learning Objectives
\begin{frame}{Learning Objectives 🎯}
    By the end of this lecture, you will:

    \begin{enumerate}
        \item Understand what tokenization is and why it matters
        \item Explain the limitations of word-level tokenization
        \item Describe how subword tokenization works (BPE, WordPiece, SentencePiece)
        \item Compare different tokenization methods and their trade-offs
        \item Implement and experiment with different tokenizers using HuggingFace
    \end{enumerate}

    \vspace{1em}
    \textbf{Central question:} How do we break text into meaningful units for AI models? 🔤
\end{frame}

% Introduction
\begin{frame}{What is Tokenization? 🔍}
    \textbf{Definition:} Converting text into smaller units (tokens)

    \vspace{0.5em}
    \textbf{Why tokenize?}
    \begin{itemize}
        \item Neural networks process numbers, not text
        \item Need discrete units to create vocabulary
        \item First step in any NLP pipeline
    \end{itemize}

    \vspace{0.5em}
    \textbf{What can tokens be?}
    \begin{itemize}
        \item \textbf{Characters:} a, b, c, ... (very fine-grained)
        \item \textbf{Words:} "hello", "world" (intuitive but limited)
        \item \textbf{Subwords:} "un", "happiness" (sweet spot! 🎯)
        \item \textbf{Bytes:} 0-255 (most fine-grained)
    \end{itemize}

    \vspace{0.5em}
    \begin{center}
        \textit{"The granularity of tokenization determines what a model can learn"}
    \end{center}
\end{frame}

\begin{frame}{The Tokenization Spectrum 📊}
    \begin{center}
        \begin{tikzpicture}[scale=0.9]
            % Spectrum line
            \draw[thick, ->, >=stealth] (0,0) -- (12,0);
            \node[below] at (0,-0.3) {Fine-grained};
            \node[below] at (12,-0.3) {Coarse-grained};

            % Character level
            \node[rectangle, draw, fill=purple!30, minimum width=2cm, minimum height=0.8cm] at (1.5,0.5) {Character};
            \node[below, font=\tiny] at (1.5,-0.8) {Vocab: 100s};
            \node[below, font=\tiny] at (1.5,-1.2) {Seq: Very long};

            % Subword level
            \node[rectangle, draw, fill=orange!30, minimum width=2cm, minimum height=0.8cm] at (6,0.5) {Subword};
            \node[below, font=\tiny] at (6,-0.8) {Vocab: 30k-50k};
            \node[below, font=\tiny] at (6,-1.2) {Seq: Medium};

            % Word level
            \node[rectangle, draw, fill=lightblue!30, minimum width=2cm, minimum height=0.8cm] at (10.5,0.5) {Word};
            \node[below, font=\tiny] at (10.5,-0.8) {Vocab: Millions};
            \node[below, font=\tiny] at (10.5,-1.2) {Seq: Short};

            % Sweet spot marker
            \draw[dashed, thick, color=darkgreen] (6,-1.5) -- (6,1.5);
            \node[above, color=darkgreen, font=\small\bfseries] at (6,1.8) {Sweet Spot! ✨};
        \end{tikzpicture}
    \end{center}

    \vspace{0.5em}
    \textbf{Key trade-off:} Vocabulary size vs. sequence length
\end{frame}

% Word-level limitations
\begin{frame}{Why Not Just Split on Spaces? 🤔}
    \textbf{The naive approach:}
    \begin{quote}
        \texttt{text.split(" ")} → Done! 🎉
    \end{quote}

    \vspace{0.5em}
    \textbf{Problem 1: Vocabulary explosion 💥}
    \begin{itemize}
        \item English: ~170,000 words in active use
        \item Each inflection counts separately: "run", "runs", "running", "ran"
        \item Compounds: "ice cream", "icecream", "ice-cream"
    \end{itemize}

    \vspace{0.5em}
    \textbf{Problem 2: Unknown words ❓}
    \begin{itemize}
        \item Out-of-vocabulary (OOV): "supercalifragilisticexpialidocious"
        \item New words: "COVID-19", "selfie", "blockchain"
        \item Typos and misspellings: "teh" instead of "the"
    \end{itemize}

    \vspace{0.5em}
    \textbf{Problem 3: No spaces in some languages! 🌏}
    \begin{itemize}
        \item Chinese: 今天天气很好 (no spaces between words)
        \item Japanese: ありがとうございます
    \end{itemize}
\end{frame}

\begin{frame}{Word-level Limitations Illustrated 📉}
    \begin{center}
        \begin{tikzpicture}
            % Vocabulary size growth
            \draw[->, thick] (0,0) -- (10,0) node[right] {Corpus size};
            \draw[->, thick] (0,0) -- (0,5) node[above] {Vocabulary};

            % Word-level curve (steep)
            \draw[thick, color=lightblue, domain=0:9] plot (\x, {0.3*\x + 0.5});
            \node[color=lightblue, font=\small] at (7,4) {Word-level 📈};

            % Subword curve (logarithmic)
            \draw[thick, color=orange, domain=0:9] plot (\x, {1.5*ln(\x+1)});
            \node[color=orange, font=\small] at (7,2.5) {Subword-level 📊};

            % Character curve (flat)
            \draw[thick, color=purple, domain=0:9] plot (\x, {0.5});
            \node[color=purple, font=\small] at (7,0.8) {Character-level →};

            % Annotations
            \node[font=\tiny] at (9,-0.5) {More data};
        \end{tikzpicture}
    \end{center}

    \vspace{0.5em}
    \textbf{Observation:} Word-level vocabulary keeps growing!

    Subword vocabulary stabilizes at reasonable size.
\end{frame}

\begin{frame}[fragile]{The OOV Problem in Action 🚨}
    \begin{lstlisting}[language=Python]
# Simple word-level vocabulary
vocab = {"hello", "world", "the", "cat", "sat"}

def tokenize_word_level(text, vocab):
    tokens = text.lower().split()
    result = []
    for token in tokens:
        if token in vocab:
            result.append(token)
        else:
            result.append("<UNK>")  # Unknown token
    return result

# Example
text = "Hello! The cat jumped"
tokens = tokenize_word_level(text, vocab)
print(tokens)
# Output: ['hello', 'the', 'cat', '<UNK>']
#                               ^^^ "jumped" is unknown!
    \end{lstlisting}

    \textbf{Problem:} We lose information! "jumped" becomes meaningless \texttt{<UNK>}
\end{frame}

% Subword Tokenization Introduction
\begin{frame}{The Solution: Subword Tokenization 🎯}
    \textbf{Key insight:} Most words are made of smaller meaningful pieces!

    \vspace{0.5em}
    \textbf{Examples:}
    \begin{itemize}
        \item "unhappiness" = "un" + "happiness"
        \item "preprocessing" = "pre" + "process" + "ing"
        \item "antiestablishmentarianism" = "anti" + "establish" + "ment" + "arian" + "ism"
    \end{itemize}

    \vspace{0.5em}
    \textbf{Benefits:}
    \begin{itemize}
        \item ✅ Fixed, manageable vocabulary (30k-50k tokens)
        \item ✅ Handle unknown words by breaking into known parts
        \item ✅ Capture morphology: prefixes, suffixes, stems
        \item ✅ Language-agnostic: works for any language!
    \end{itemize}

    \vspace{0.5em}
    \textbf{How?} Learn common character sequences from data! 📊
\end{frame}

\begin{frame}{Visual: How "unhappiness" Gets Tokenized 📊}
    \begin{center}
        \begin{tikzpicture}[scale=0.9]
            % Word-level
            \node[rectangle, draw, fill=lightblue!30, minimum width=3cm, minimum height=1cm] at (0,3.5) {\texttt{unhappiness}};
            \node[anchor=east] at (-0.5,3.5) {Word-level:};
            \node[anchor=west, font=\small] at (2,3.5) {1 token};

            % Subword BPE
            \node[rectangle, draw, fill=orange!30, minimum width=1.2cm, minimum height=1cm] at (-0.9,2) {\texttt{un}};
            \node[rectangle, draw, fill=orange!30, minimum width=2cm, minimum height=1cm] at (0.7,2) {\texttt{happiness}};
            \node[anchor=east] at (-2.5,2) {Subword (BPE):};
            \node[anchor=west, font=\small] at (2.2,2) {2 tokens};

            % Alternative BPE
            \node[rectangle, draw, fill=orange!20, minimum width=1.2cm, minimum height=0.8cm, font=\small] at (-1.5,0.8) {\texttt{un}};
            \node[rectangle, draw, fill=orange!20, minimum width=1cm, minimum height=0.8cm, font=\small] at (-0.3,0.8) {\texttt{happ}};
            \node[rectangle, draw, fill=orange!20, minimum width=0.8cm, minimum height=0.8cm, font=\small] at (0.7,0.8) {\texttt{i}};
            \node[rectangle, draw, fill=orange!20, minimum width=1cm, minimum height=0.8cm, font=\small] at (1.7,0.8) {\texttt{ness}};
            \node[anchor=east, font=\small] at (-2.5,0.8) {Alternative:};
            \node[anchor=west, font=\small] at (2.5,0.8) {4 tokens};

            % Character-level
            \node[rectangle, draw, fill=darkgreen!30, minimum width=0.4cm, minimum height=0.8cm] at (-3.3,-0.5) {\tiny \texttt{u}};
            \node[rectangle, draw, fill=darkgreen!30, minimum width=0.4cm, minimum height=0.8cm] at (-2.9,-0.5) {\tiny \texttt{n}};
            \node[rectangle, draw, fill=darkgreen!30, minimum width=0.4cm, minimum height=0.8cm] at (-2.5,-0.5) {\tiny \texttt{h}};
            \node[rectangle, draw, fill=darkgreen!30, minimum width=0.4cm, minimum height=0.8cm] at (-2.1,-0.5) {\tiny \texttt{a}};
            \node[rectangle, draw, fill=darkgreen!30, minimum width=0.4cm, minimum height=0.8cm] at (-1.7,-0.5) {\tiny \texttt{p}};
            \node[rectangle, draw, fill=darkgreen!30, minimum width=0.4cm, minimum height=0.8cm] at (-1.3,-0.5) {\tiny \texttt{p}};
            \node[rectangle, draw, fill=darkgreen!30, minimum width=0.4cm, minimum height=0.8cm] at (-0.9,-0.5) {\tiny \texttt{i}};
            \node[rectangle, draw, fill=darkgreen!30, minimum width=0.4cm, minimum height=0.8cm] at (-0.5,-0.5) {\tiny \texttt{n}};
            \node[rectangle, draw, fill=darkgreen!30, minimum width=0.4cm, minimum height=0.8cm] at (-0.1,-0.5) {\tiny \texttt{e}};
            \node[rectangle, draw, fill=darkgreen!30, minimum width=0.4cm, minimum height=0.8cm] at (0.3,-0.5) {\tiny \texttt{s}};
            \node[rectangle, draw, fill=darkgreen!30, minimum width=0.4cm, minimum height=0.8cm] at (0.7,-0.5) {\tiny \texttt{s}};
            \node[anchor=east] at (-4.5,-0.5) {Character-level:};
            \node[anchor=west, font=\small] at (1.2,-0.5) {11 tokens};
        \end{tikzpicture}
    \end{center}

    \vspace{0.5em}
    \textbf{Trade-off:} More tokens = longer sequences vs. smaller vocabulary
\end{frame}

% Byte-Pair Encoding
\begin{frame}{Byte-Pair Encoding (BPE) 🔧}
    \textbf{Original use:} Data compression (1994)

    \textbf{Adapted for NLP:} Sennrich et al. (2016)

    \vspace{0.5em}
    \textbf{Core idea:} Iteratively merge the most frequent pair of characters/subwords

    \vspace{0.5em}
    \textbf{Algorithm:}
    \begin{enumerate}
        \item Start with character-level vocabulary
        \item Count all adjacent pairs in corpus
        \item Merge most frequent pair → create new token
        \item Repeat until desired vocabulary size
    \end{enumerate}

    \vspace{0.5em}
    \textbf{Used by:} GPT, GPT-2, GPT-3, RoBERTa, BART, many others!

    \vspace{0.5em}
    📖 \textbf{Reference:} Sennrich, Haddow, \& Birch (2016). Neural Machine Translation of Rare Words with Subword Units. \textit{ACL}.
\end{frame}

\begin{frame}{BPE Example: Step by Step 👣}
    \textbf{Training data:} "low low low lower lower newest newest newest newest widest"

    \vspace{0.5em}
    \textbf{Initial tokens (characters):}
    \begin{center}
        \texttt{l, o, w, e, r, n, s, t, i, d}
    \end{center}

    \vspace{0.5em}
    \textbf{Iteration 1:} Most frequent pair = \texttt{e, s} (appears 4 times)
    \begin{itemize}
        \item Merge → new token: \texttt{es}
        \item Vocabulary: \texttt{l, o, w, e, r, n, s, t, i, d, es}
    \end{itemize}

    \vspace{0.3em}
    \textbf{Iteration 2:} Most frequent = \texttt{es, t}
    \begin{itemize}
        \item Merge → new token: \texttt{est}
        \item Vocabulary: \texttt{l, o, w, e, r, n, s, t, i, d, es, est}
    \end{itemize}

    \vspace{0.3em}
    \textbf{Iteration 3:} Most frequent = \texttt{l, o}
    \begin{itemize}
        \item Merge → new token: \texttt{lo}
    \end{itemize}

    \vspace{0.3em}
    Continue until reaching target vocabulary size (e.g., 30,000)...
\end{frame}

\begin{frame}{BPE Visualization 🎨}
    \begin{center}
        \begin{tikzpicture}[
            token/.style={rectangle, draw, fill=lightblue!30, minimum width=0.6cm, minimum height=0.6cm, font=\tiny},
            merge/.style={rectangle, draw, fill=orange!30, minimum width=1.3cm, minimum height=0.6cm, font=\tiny}
        ]
            % Step 0: Characters
            \node[font=\small] at (-1,3) {Step 0:};
            \node[token] at (0,3) {l};
            \node[token] at (0.7,3) {o};
            \node[token] at (1.4,3) {w};
            \node[token] at (2.1,3) {e};
            \node[token] at (2.8,3) {s};
            \node[token] at (3.5,3) {t};

            % Step 1: Merge e+s
            \node[font=\small] at (-1,2) {Step 1:};
            \node[token] at (0,2) {l};
            \node[token] at (0.7,2) {o};
            \node[token] at (1.4,2) {w};
            \node[merge] at (2.65,2) {es};
            \node[token] at (3.5,2) {t};

            % Step 2: Merge es+t
            \node[font=\small] at (-1,1) {Step 2:};
            \node[token] at (0,1) {l};
            \node[token] at (0.7,1) {o};
            \node[token] at (1.4,1) {w};
            \node[merge, fill=darkgreen!30, minimum width=2cm] at (3,1) {est};

            % Step 3: Merge l+o
            \node[font=\small] at (-1,0) {Step 3:};
            \node[merge] at (0.5,0) {lo};
            \node[token] at (1.4,0) {w};
            \node[merge, fill=darkgreen!30, minimum width=2cm] at (3,0) {est};

            % Arrows
            \draw[->, thick, color=orange] (2.45,2.8) -- (2.65,2.2);
            \draw[->, thick, color=darkgreen] (2.95,1.8) -- (3,1.2);
            \draw[->, thick, color=orange] (0.35,2.8) -- (0.5,0.2);
        \end{tikzpicture}
    \end{center}

    \vspace{0.5em}
    \textbf{Result:} "lowest" → \texttt{[lo, w, est]}

    Captures common patterns without explicit linguistic rules!
\end{frame}

\begin{frame}[fragile]{BPE in Practice with HuggingFace 🤗}
    \begin{lstlisting}[language=Python]
from transformers import AutoTokenizer

# GPT-2 uses BPE
tokenizer = AutoTokenizer.from_pretrained("gpt2")

text = "I'm learning about tokenization!"

# Tokenize
tokens = tokenizer.tokenize(text)
print("Tokens:", tokens)
# Output: ['I', "'m", 'Ġlearning', 'Ġabout', 'Ġtoken', 'ization', '!']
# Note: 'Ġ' represents space

# Get token IDs
token_ids = tokenizer.encode(text)
print("Token IDs:", token_ids)
# Output: [40, 1101, 4673, 546, 11241, 1634, 0]

# Decode back
decoded = tokenizer.decode(token_ids)
print("Decoded:", decoded)
# Output: "I'm learning about tokenization!"
    \end{lstlisting}
\end{frame}

% WordPiece
\begin{frame}{WordPiece Tokenization 🧩}
    \textbf{Similar to BPE, but with a twist!}

    \vspace{0.5em}
    \textbf{Key difference:} Instead of merging most \textit{frequent} pair, merge pair that maximizes \textit{likelihood} of training data

    \vspace{0.5em}
    \textbf{Merge criterion:}
    \[
    \text{score}(x, y) = \frac{\text{freq}(xy)}{\text{freq}(x) \times \text{freq}(y)}
    \]

    \vspace{0.5em}
    \textbf{Intuition:} Prefer merging pairs that are "surprisingly common" together

    \vspace{0.5em}
    \textbf{Used by:} BERT, DistilBERT, Electra

    \vspace{0.5em}
    \textbf{Special tokens:}
    \begin{itemize}
        \item \texttt{\#\#} prefix for continuation subwords
        \item Example: "playing" → \texttt{["play", "\#\#ing"]}
    \end{itemize}

    \vspace{0.5em}
    📖 \textbf{Reference:} Wu et al. (2016). Google's Neural Machine Translation System. \textit{arXiv}.
\end{frame}

\begin{frame}[fragile]{WordPiece Example with BERT 🔍}
    \begin{lstlisting}[language=Python]
from transformers import AutoTokenizer

# BERT uses WordPiece
tokenizer = AutoTokenizer.from_pretrained("bert-base-uncased")

text = "I'm learning about tokenization!"

tokens = tokenizer.tokenize(text)
print("Tokens:", tokens)
# Output: ['i', "'", 'm', 'learning', 'about', 'token', '##ization', '!']
#                                                        ^^^ Note the ##

# Compare with longer word
text2 = "unhappiness"
tokens2 = tokenizer.tokenize(text2)
print("Tokens:", tokens2)
# Output: ['un', '##hap', '##pin', '##ess']
# Breaks into morphological components!
    \end{lstlisting}

    \vspace{0.5em}
    \textbf{Notice:} BERT lowercases by default (unless using cased model)
\end{frame}

% SentencePiece
\begin{frame}{SentencePiece: Language-Agnostic Tokenization 🌍}
    \textbf{Problem with BPE/WordPiece:}
    \begin{itemize}
        \item Assume text is pre-tokenized (words separated by spaces)
        \item Doesn't work for languages without spaces (Chinese, Japanese, Thai...)
    \end{itemize}

    \vspace{0.5em}
    \textbf{SentencePiece solution:}
    \begin{itemize}
        \item Treats input as raw character stream (no pre-tokenization!)
        \item Spaces are just another character: \texttt{▁}
        \item Truly language-agnostic
    \end{itemize}

    \vspace{0.5em}
    \textbf{Two algorithms:}
    \begin{enumerate}
        \item \textbf{Unigram Language Model:} Start with large vocabulary, iteratively remove tokens
        \item \textbf{BPE:} Same as before, but on raw text
    \end{enumerate}

    \vspace{0.5em}
    \textbf{Used by:} T5, ALBERT, XLNet, mT5 (multilingual models)

    \vspace{0.5em}
    📖 \textbf{Reference:} Kudo \& Richardson (2018). SentencePiece: A simple and language independent approach to subword tokenization. \textit{EMNLP}.
\end{frame}

\begin{frame}[fragile]{SentencePiece in Action 🚀}
    \begin{lstlisting}[language=Python]
from transformers import AutoTokenizer

# T5 uses SentencePiece
tokenizer = AutoTokenizer.from_pretrained("t5-small")

text = "I'm learning about tokenization!"

tokens = tokenizer.tokenize(text)
print("Tokens:", tokens)
# Output: ['▁I', "'", 'm', '▁learning', '▁about', '▁token', 'ization', '!']
#         ^^^ Note the ▁ for spaces

# Works seamlessly with other languages!
text_chinese = "今天天气很好"
tokens_cn = tokenizer.tokenize(text_chinese)
print("Chinese:", tokens_cn)
# Breaks into characters/subwords without needing pre-tokenization
    \end{lstlisting}

    \vspace{0.5em}
    \textbf{Key advantage:} No language-specific preprocessing required! 🌏
\end{frame}

% Comparison
\begin{frame}{Tokenization Methods Comparison 📋}
    \begin{table}
        \centering
        \small
        \begin{tabular}{lllp{3.5cm}}
            \toprule
            \textbf{Method} & \textbf{Approach} & \textbf{Used By} & \textbf{Key Feature} \\
            \midrule
            BPE & Frequency-based & GPT-2, GPT-3, & Simple, effective, \\
            & merging & RoBERTa & bottom-up \\
            \midrule
            WordPiece & Likelihood-based & BERT, & Statistically \\
            & merging & DistilBERT & principled \\
            \midrule
            SentencePiece & No pre-tokenization, & T5, ALBERT, & Language-agnostic, \\
            (Unigram) & top-down pruning & mT5 & multilingual \\
            \midrule
            SentencePiece & No pre-tokenization, & XLNet, & Combines BPE \\
            (BPE) & bottom-up & mBART & with raw text \\
            \bottomrule
        \end{tabular}
    \end{table}

    \vspace{0.5em}
    📖 \textbf{HuggingFace Resources:}
    \begin{itemize}
        \item \href{https://huggingface.co/learn/nlp-course/chapter2/4}{Chapter 2.4: Behind the Pipeline - Tokenizers}
        \item \href{https://huggingface.co/learn/nlp-course/chapter6}{Chapter 6: Tokenizers (full chapter)}
    \end{itemize}
\end{frame}

\begin{frame}[fragile]{Comparing Tokenizers Side-by-Side 🔬}
    \begin{lstlisting}[language=Python]
from transformers import AutoTokenizer

text = "The unhappiest researchers couldn't preprocess data!"

models = {
    "GPT-2 (BPE)": "gpt2",
    "BERT (WordPiece)": "bert-base-uncased",
    "T5 (SentencePiece)": "t5-small"
}

for name, model_name in models.items():
    tokenizer = AutoTokenizer.from_pretrained(model_name)
    tokens = tokenizer.tokenize(text)
    print(f"\n{name}:")
    print(f"  Tokens ({len(tokens)}): {tokens}")

# Observe:
# - How each handles "unhappiest"
# - Different vocabulary sizes
# - Different handling of spaces
# - Number of tokens varies!
    \end{lstlisting}
\end{frame}

% Special Tokens
\begin{frame}{Special Tokens 🎟️}
    \textbf{Most tokenizers add special tokens for specific purposes:}

    \vspace{0.5em}
    \begin{table}
        \centering
        \begin{tabular}{lll}
            \toprule
            \textbf{Token} & \textbf{Purpose} & \textbf{Model} \\
            \midrule
            \texttt{[CLS]} & Classification token & BERT \\
            \texttt{[SEP]} & Separator between segments & BERT \\
            \texttt{[PAD]} & Padding to same length & Most models \\
            \texttt{[UNK]} & Unknown/rare tokens & Most models \\
            \texttt{[MASK]} & Masked token for training & BERT \\
            \texttt{<s>}, \texttt{</s>} & Start/end of sequence & GPT-2, T5 \\
            \texttt{<|endoftext|>} & Document boundary & GPT-2 \\
            \bottomrule
        \end{tabular}
    \end{table}

    \vspace{0.5em}
    \textbf{Example usage:}
    \begin{center}
        \texttt{[CLS] The cat sat [SEP] The dog ran [SEP]}
    \end{center}
\end{frame}

\begin{frame}[fragile]{Working with Special Tokens 🔧}
    \begin{lstlisting}[language=Python]
from transformers import AutoTokenizer

tokenizer = AutoTokenizer.from_pretrained("bert-base-uncased")

# Special tokens are automatically added
text = "Hello world"
encoded = tokenizer(text, return_tensors="pt")
print("Input IDs:", encoded['input_ids'])
# [101, 7592, 2088, 102]
#  ^^^             ^^^ [CLS] and [SEP] added!

# Decode with/without special tokens
full_decode = tokenizer.decode(encoded['input_ids'][0])
print("With special:", full_decode)
# Output: "[CLS] hello world [SEP]"

clean_decode = tokenizer.decode(encoded['input_ids'][0],
                                skip_special_tokens=True)
print("Without special:", clean_decode)
# Output: "hello world"
    \end{lstlisting}
\end{frame}

% Practical Considerations
\begin{frame}{Vocabulary Size Trade-offs 🎚️}
    \textbf{Smaller vocabulary (e.g., 10k tokens):}
    \begin{itemize}
        \item ➕ Faster training (smaller embedding matrix)
        \item ➕ Less memory
        \item ➖ Longer sequences (more subwords per word)
        \item ➖ May lose semantic information
    \end{itemize}

    \vspace{0.5em}
    \textbf{Larger vocabulary (e.g., 100k tokens):}
    \begin{itemize}
        \item ➕ Shorter sequences (closer to word-level)
        \item ➕ Better semantic preservation
        \item ➖ Slower training (larger embeddings)
        \item ➖ More memory required
    \end{itemize}

    \vspace{0.5em}
    \textbf{Sweet spot:} 30k-50k tokens for most models

    \begin{center}
        GPT-2: 50k | BERT: 30k | T5: 32k
    \end{center}
\end{frame}

\begin{frame}{Tokenization Pitfalls and Gotchas ⚠️}
    \textbf{Common issues to watch out for:}

    \begin{enumerate}
        \item \textbf{Tokenizer-model mismatch}
        \begin{itemize}
            \item Always use the tokenizer that matches your model!
            \item GPT-2 tokenizer ≠ BERT tokenizer
        \end{itemize}

        \item \textbf{Maximum sequence length}
        \begin{itemize}
            \item BERT: 512 tokens | GPT-2: 1024 | GPT-3: 2048
            \item Text gets truncated if too long!
        \end{itemize}

        \item \textbf{Case sensitivity}
        \begin{itemize}
            \item \texttt{bert-base-uncased} lowercases everything
            \item \texttt{bert-base-cased} preserves case
        \end{itemize}

        \item \textbf{Rare words → many tokens}
        \begin{itemize}
            \item "antidisestablishmentarianism" → 10+ tokens
            \item Can hit sequence limit faster than expected!
        \end{itemize}

        \item \textbf{Special characters}
        \begin{itemize}
            \item Emoji, Unicode, accents may be split unexpectedly
        \end{itemize}
    \end{enumerate}
\end{frame}

\begin{frame}[fragile]{Debugging Tokenization 🐛}
    \begin{lstlisting}[language=Python]
from transformers import AutoTokenizer

tokenizer = AutoTokenizer.from_pretrained("gpt2")

text = "The antidisestablishmentarianism debate continues! 😊"

# Tokenize and see details
tokens = tokenizer.tokenize(text)
print(f"Tokens ({len(tokens)}): {tokens}\n")

# Get IDs and decode
token_ids = tokenizer.encode(text)
for i, (token, token_id) in enumerate(zip(tokens, token_ids)):
    decoded = tokenizer.decode([token_id])
    print(f"{i:2d}. ID {token_id:5d} | Token: {token:20s} | Decoded: {decoded}")

# Check vocabulary size
print(f"\nVocabulary size: {tokenizer.vocab_size}")

# Check special tokens
print(f"Special tokens: {tokenizer.all_special_tokens}")
    \end{lstlisting}
\end{frame}

% Connection to Human Language
\begin{frame}{Connection to Human Language Learning 👶}
    \textbf{Reminder from infant learning (Saffran et al., 1996):}
    \begin{itemize}
        \item Babies track statistical regularities in speech
        \item Identify word boundaries from transitional probabilities
        \item No explicit rules—just patterns from exposure!
    \end{itemize}

    \vspace{0.5em}
    \textbf{Parallel with subword tokenization:}
    \begin{itemize}
        \item BPE: Merge frequent character pairs → discover common morphemes
        \item Infants: Track frequent syllable pairs → discover words
        \item Both are \textit{data-driven}, not rule-based!
    \end{itemize}

    \vspace{0.5em}
    \textbf{Key difference:}
    \begin{itemize}
        \item Infants: Online, real-time, multimodal (sound + context)
        \item BPE: Batch, text-only, no grounding
    \end{itemize}

    \vspace{0.5em}
    📖 \textbf{Reference:} Saffran, Aslin, \& Newport (1996). Statistical learning by 8-month-old infants. \textit{Science}.
\end{frame}

\begin{frame}{Statistical Learning in Action 📊}
    \begin{columns}[T]
        \begin{column}{0.48\textwidth}
            \textbf{Infant Learning:}

            Input stream:
            \begin{center}
                \texttt{bidakupadotigolabu...}
            \end{center}

            Learn high-probability sequences:
            \begin{itemize}
                \item \texttt{bi-da-ku} (word)
                \item \texttt{pa-do-ti} (word)
            \end{itemize}

            Detect low-probability boundaries:
            \begin{itemize}
                \item \texttt{ku | pa} (boundary)
            \end{itemize}
        \end{column}
        \begin{column}{0.48\textwidth}
            \textbf{BPE Learning:}

            Input corpus:
            \begin{center}
                \texttt{low low lower...}
            \end{center}

            Merge high-frequency pairs:
            \begin{itemize}
                \item \texttt{l+o → lo}
                \item \texttt{lo+w → low}
            \end{itemize}

            Build vocabulary:
            \begin{itemize}
                \item \texttt{low, lower, lowest}
            \end{itemize}
        \end{column}
    \end{columns}

    \vspace{1em}
    \textbf{Both use distributional statistics to discover structure!} 🎯
\end{frame}

% Practical Exercise
\begin{frame}{Hands-On Exercise 🧪}
    \textbf{Experiment with different tokenizers!}

    \begin{enumerate}
        \item Choose 3 models: GPT-2, BERT, T5
        \item Test on diverse texts:
        \begin{itemize}
            \item Standard English: "The cat sat on the mat"
            \item Complex words: "antidisestablishmentarianism"
            \item Contractions: "I'm, you're, won't"
            \item Typos: "teh qiuck brown fox"
            \item Emoji: "I love this! 😊🎉"
            \item Other languages: "今天天气很好" (Chinese)
        \end{itemize}
        \item Compare results:
        \begin{itemize}
            \item Number of tokens
            \item How words are split
            \item Handling of unknown/rare words
        \end{itemize}
        \item Reflect:
        \begin{itemize}
            \item Which tokenizer works best for your use case?
            \item What are the trade-offs?
        \end{itemize}
    \end{enumerate}
\end{frame}

\begin{frame}{Discussion Questions 💭}
    \begin{enumerate}
        \item \textbf{Linguistics vs. Statistics:}
        \begin{itemize}
            \item BPE discovers morphemes (un-, -ing, -ness) without linguistic rules. Is this "learning" morphology, or just pattern matching?
        \end{itemize}

        \item \textbf{Cross-lingual tokenization:}
        \begin{itemize}
            \item Should we use the same tokenizer for all languages? What are the trade-offs?
        \end{itemize}

        \item \textbf{Semantic preservation:}
        \begin{itemize}
            \item Does breaking "unhappy" into ["un", "happy"] preserve meaning? What about "butterfly"?
        \end{itemize}

        \item \textbf{Human vs. machine:}
        \begin{itemize}
            \item Humans don't consciously tokenize words. Why do machines need to?
        \end{itemize}

        \item \textbf{Future directions:}
        \begin{itemize}
            \item Will we move toward character-level or byte-level models that don't need tokenization?
        \end{itemize}
    \end{enumerate}
\end{frame}

% Resources and Summary
\begin{frame}{Primary References 📚}
    \textbf{Foundational papers:}
    \begin{itemize}
        \item Sennrich, Haddow, \& Birch (2016). Neural Machine Translation of Rare Words with Subword Units. \textit{ACL}.
        \begin{itemize}
            \item Introduced BPE for NLP
        \end{itemize}

        \item Kudo \& Richardson (2018). SentencePiece: A simple and language independent approach. \textit{EMNLP}.
        \begin{itemize}
            \item Language-agnostic tokenization
        \end{itemize}

        \item Wu et al. (2016). Google's Neural Machine Translation System. \textit{arXiv}.
        \begin{itemize}
            \item WordPiece algorithm
        \end{itemize}
    \end{itemize}

    \vspace{0.5em}
    \textbf{HuggingFace resources:}
    \begin{itemize}
        \item \href{https://huggingface.co/learn/nlp-course/chapter2/4}{Chapter 2.4: Tokenizers}
        \item \href{https://huggingface.co/learn/nlp-course/chapter6}{Chapter 6: Tokenizers (detailed)}
    \end{itemize}
\end{frame}

\begin{frame}{Key Takeaways 🎯}
    \begin{enumerate}
        \item \textbf{Tokenization is fundamental}
        \begin{itemize}
            \item Bridges raw text and neural networks
        \end{itemize}

        \item \textbf{Word-level has major limitations}
        \begin{itemize}
            \item OOV problem, vocabulary explosion, language-dependent
        \end{itemize}

        \item \textbf{Subword tokenization is the sweet spot}
        \begin{itemize}
            \item Balanced vocabulary size and sequence length
        \end{itemize}

        \item \textbf{Different methods, similar principles}
        \begin{itemize}
            \item BPE: frequency-based | WordPiece: likelihood | SentencePiece: language-agnostic
        \end{itemize}

        \item \textbf{Statistical learning connects humans and machines}
        \begin{itemize}
            \item Both discover structure from distributional patterns
        \end{itemize}

        \item \textbf{Always match tokenizer to model!}
        \begin{itemize}
            \item Critical for correct predictions
        \end{itemize}
    \end{enumerate}
\end{frame}

\begin{frame}{Looking Ahead 🔮}
    \textbf{Next lecture:} POS Tagging \& Sentiment Analysis

    \vspace{0.5em}
    \textbf{How tokenization connects:}
    \begin{itemize}
        \item POS tagging: Token-level classification
        \item Sentiment: Sequence-level classification
        \item Both depend on good tokenization!
    \end{itemize}

    \vspace{1em}
    \textbf{Prepare by:}
    \begin{itemize}
        \item Experimenting with HuggingFace tokenizers
        \item Thinking about: How does token granularity affect downstream tasks?
        \item Exploring: Tokenizer artifacts and their impact
    \end{itemize}
\end{frame}

\begin{frame}[standout]
    Questions? 🤔

    \vspace{2em}

    \normalsize
    Next: Lecture 6 - POS Tagging \& Sentiment Analysis
\end{frame}

\end{document}
